\chapter{Einleitung}
\label{chap:einleitung}

Moderne Technologien wie \gls{gls-KI} und die Integration in Telefonsysteme führen dazu, dass sich sowohl die Art der Interaktion als auch die Qualität und Komplexität dieser Systeme weiterentwickeln. Sprachassistenten ermöglichen die automatische Erkennung menschlicher Sprache und können beispielsweise individuelle Funktionen zur Steuerung von Smart-Home-Komponenten übernehmen. Die Integration von \gls{gls-LLM} erweitert diese Funktionen um dialogorientierte Steuerung und Kontextverarbeitung über mehrere Gesprächsbeiträge hinweg. Dadurch werden natürlichere und flexiblere Dialoge möglich, die klassische Telefonie sowie moderne Sprachplattformen ergänzen und neue Formen der Mensch-Maschine-Kommunikation schaffen \parencite{Schmiedchen2026KI}. Hierbei stellt sich die Frage, wie \gls{gls-KMU} mit begrenzter Mitarbeiter- und Umsatzgrö\ss e von solchen Technologien profitieren können.

\section{Problemstellung und Zielsetzung}
\gls{gls-VoIP}-Systeme, die auf dem \gls{gls-SIP}-Protokoll basieren, sind häufig stark an spezielle Telefonie-Hardware gebunden und bringen kaum intelligente Sprachverarbeitung mit. Die Verbindung solcher Systeme mit einem \gls{gls-LLM} eröffnet neue Möglichkeiten, etwa für automatisierte Sprachdialoge, die Kundensysteme effizient unterstützen und
gleichzeitig natürlich wirken. Forschung zur Integration von Sprachverarbeitung mit \gls{gls-LLM} zeigt, dass die Kombination von Audio-Encodern mit Sprachmodellen robuste multilinguale Spracherkennung ermöglicht \parencite{whisper_quant_2025}. Au\ss erdem wird deutlich, dass sich klassische \gls{gls-SIP}-basierte Telefonsysteme mit modernen \gls{gls-KI}-Dialogkomponenten und containerbasierten Architekturen kombinieren lassen \parencite{Antonova2025}. In der Praxis zeigt sich jedoch, dass die Übertragung und Verarbeitung von Audio zwischen virtualisierten Systemkomponenten eine gro\ss e Herausforderung darstellt. In containerbasierten Umgebungen, in denen verschiedene Dienste getrennt voneinander laufen, müssen Sprachdaten zuverlässig und mit möglichst niedriger Latenz zwischen diesen Diensten übertragen werden \parencite{ReliabilityVirtualizedNetworks}.

Neuere Studien weisen darauf hin, dass \gls{gls-SIP}-Instanzen und die dazugehörigen Media-Proxy in Container-Umgebungen eingebunden werden sollten. Ein Media-Proxy leitet Audio- und Videoströme weiter, passt sie bei Bedarf an oder transkodiert sie, damit unterschiedliche Protokolle miteinander kompatibel sind. Dabei ist es wichtig, dass die Verwaltung der Dienste automatisiert und orchestriert erfolgt, um Skalierbarkeit und Fehlertoleranz sicherzustellen \parencite{Antonova2025}. Die Realisierung solcher Systeme auf Basis von Open-Source-Softwarekomponenten und Container-Orchestrierung mit Docker ermöglicht eine reproduzierbare, kosteneffiziente und flexible Bereitstellung. Im Rahmen dieser Arbeit wurde durch eine iterative Evaluierung verschiedener \gls{gls-SIP}-Implementierungen eine praxisnahe Architektur entwickelt. Ausgehend von baresip als experimentellem Tool zur Überprüfung der SIP-Registrierung und Protokollkompatibilität wurde die Architektur schrittweise mit Asterisk als stabilem \gls{gls-SIP}-Backend und LiveKit als Integrations-Framework für \gls{gls-KI}-Sprachverarbeitungsdienste realisiert.

Das \gls{gls-RTP} übernimmt die Echtzeitübertragung von Audio und Video, während \gls{gls-WebRTC} die Kommunikation im Browser ermöglicht. Für die Medienverarbeitung ist eine RTP-Bridge erforderlich, die Mediendaten zwischen den Diensten weiterleitet, anpasst oder transkodiert. Sicherheitsfunktionen wie das \gls{gls-SRTP}, zur Verschlüsselung der Mediendaten und \gls{gls-DTLS} für die sichere Schlüsselverteilung, schützen die übertragenen Daten. Die Interoperabilität zwischen \gls{gls-SIP}/\gls{gls-RTP} und \gls{gls-WebRTC} stellt sicher, dass unterschiedliche Clients zuverlässig miteinander kommunizieren können. Ohne diese Komponenten lassen sich stabile, bidirektionale Audio-Pipelines mit kontrollierter Latenz kaum realisieren\\ \parencite{Antonova2025, ReliabilityVirtualizedNetworks}.

\section{Aufbau der Arbeit}
 Kapitel \ref{chap:grundlagen} beschreibt die technischen Grundlagen zu allen genutzten Technologien und Protokollen. Anschlie\ss end enthält Kapitel \ref{chap:anforderungen} Anforderungen für die Implementierung. Das Kapitel \ref{chap:architektur} behandelt Architekturthematiken und Kapitel \ref{chap:marktanalyse} gibt Einblicke in den Markt von \gls{gls-KI}-basierten Systemen. Im nächsten Kapitel \ref{chap:implementierung} werden die prototypische Umsetzung und die Integration der zentralen Komponenten dokumentiert. Danach präsentiert Kapitel \ref{chap:evaluation} die Messergebnisse und deren Auswertung.
Das abschlie\ss ende Kapitel \ref{chap:diskussion} vereint die Interpretation der Ergebnisse mit den zentralen Erkenntnissen.

\section{Beitrag der Arbeit}
Der Beitrag der Arbeit umfasst: (1) eine Architektur für die Kombination von \gls{gls-SIP}-/\gls{gls-VoIP}-Systemen mit einem \gls{gls-LLM}-Dienst in virtualisierten Containern, (2) eine prototypische Implementierung der genutzten Komponenten mit reproduzierbarem Setup, (3) eine Evaluation wesentlicher Qualitätsmerkmale und (4) einen praxisorientierten Leitfaden mit technischen Handlungsempfehlungen zur Nutzung des prototypischen Systems. \\

\section{Methodik}

Bei der Realisierung dieser Arbeit wurde ein iterativer und experimenteller Ansatz verfolgt. Der Fokus lag auf dem Nachweis der technischen Machbarkeit der Integration von \gls{gls-SIP}-Systemen mit \gls{gls-KI}-Sprachverarbeitung. Baresip\footnote{Offizielle Webseite: \url{https://github.com/baresip/baresip}} diente als experimentelles Tool zur Überprüfung grundlegender \gls{gls-SIP}-Registrierungs- und Protokoll-Kompatibilitätsfragen. Aufgrund der begrenzten Dialplan-Logik wurde zügig auf Asterisk gewechselt, um die grundlegenden Mechanismen einer \gls{gls-SIP}-Trunk-Registrierung in einer produktionsnahen Umgebung abzubilden. Diese Phase führte zur Auswahl von Asterisk als stabiles und produktiv einsetzbares \gls{gls-SIP}-Backend. Das LiveKit-Framework\footnote{Offizielle Webseite: \url{https://livekit.io/}} wurde anschlie\ss end als geeignete Lösung zur Integration mit \gls{gls-KI}-Sprachverarbeitungs-\\diensten implementiert. Da der Schwerpunkt auf der finalen, funktionsfähigen Architektur liegt, werden die Evaluierungsphasen der einzelnen Komponenten nur so weit dokumentiert, wie sie für das Verständnis des Gesamtsystems notwendig sind. Außerdem wurde eine \gls{gls-UI} als praktisches Steuerungstool für reine \gls{gls-WebRTC} Telefonie ohne \gls{gls-SIP} Integration entwickelt, um die grundlegenden \gls{gls-KI}-Komponenten zu prüfen. Die wissenschaftliche Evaluation konzentriert sich jedoch auf die \gls{gls-SIP}-basierte Telefonie-Integration mit \gls{gls-KI}-Sprachverarbeitung, nicht auf die \gls{gls-UI}.

Folgende Hypothesen werden überprüft und im Evaluationskapitel \ref{chap:evaluation} anhand systematischer Tests bestätigt oder verworfen.

\paragraph{H1 (Funktionalität):} Es wird angenommen, dass die bidirektionale Audioübertragung zwischen einem \gls{gls-SIP}-Telefonsystem und einem extern gehosteten \gls{gls-KI}-\\Sprachverarbeitungsdienst in einer containerbasierten Umgebung stabil aufrechterhalten werden kann, ohne dabei systematische Verbindungsabbrüche, Codec-\\Inkompatibilitäten oder Medienpfadfehler aufzuweisen.\\

\textit{Kontext:} Nach RFC~3261 regelt \gls{gls-SIP} die Signalisierung und Codec-Verhandlung in \gls{gls-VoIP}-Systemen. Die Standards RFC~5245 (Interactive Connectivity Establishment), RFC~5389 (Session Traversal Utilities for NAT) und RFC~5766 (Traversal Using Relays around NAT) beschreiben Mechanismen zur Überwindung von Netzwerk-Adressumsetzung und Firewall-Hindernissen, die für die Herstellung stabiler Audio-Kommunikationspfade essentiell sind \parencite{rfc3261,rfc5245,rfc5389,rfc5766}.

\paragraph{H2 (Latenz):} Es wird angenommen, dass die in der Evaluation gemessenen Latenzbeiträge der Sprachverarbeitung (\gls{gls-STT}‑Latenz, Transkriptverzögerung, \gls{gls-TTS}‑Latenz) im praxistauglichen Bereich liegen und interaktive Kommunikation unterstützen.

\textit{Kontext:} In der klassischen Telefonie nach ITU‑T G.114 werden One‑Way‑Latenzen bis 150 ms als akzeptabel, bis 400 ms als noch tolerierbar und oberhalb von 400 ms als kritisch eingestuft \parencite{ITU-T-G114}. In dieser Arbeit werden die Latenzbeiträge der Sprachverarbeitung separat betrachtet, da sie die zentralen, systemseitig beeinflussbaren Schritte abbilden (\gls{gls-STT}, Transkriptverzögerung in der \gls{gls-SIP}-Pipeline, \gls{gls-TTS}). Externe Anteile (Routing, Netz, \gls{gls-SIP}‑Signalisierung) wirken zusätzlich, werden hier aber nicht gemessen.

\paragraph{H3 (Ressourceneffizienz):} Es wird angenommen, dass die eingesetzte containerbasierte Architektur auf Basis von Standard-Serverhardware betrieben werden kann und dabei Systemressourcen (\gls{gls-CPU}, \gls{gls-RAM}, Speicherbandbreite) effizient ausnutzt, sodass keine übermä\ss ige Überprovisionierung erforderlich ist.

\textit{Kontext:} Nach den \gls{gls-ETSI}-\gls{gls-NFV}-Richtlinien besteht das Ziel der Netzwerk-\\Funktionsvirtualisierung darin, spezialisierte Hardware durch software- und containerbasierte Dienste zu ersetzen, die auf Standard-Servern betrieben werden können. Praktische Studien zu solchen Kommunikationssystemen zeigen, dass mit optimierter Dimensionierung stabile und performante Deployments auch auf ressourcenlimitierter Hardware möglich sind \parencite{ETSI-NFV,ETSI-REL,Antonova2025}.

\paragraph{H4 (Skalierbarkeit):} Es wird angenommen, dass die containerbasierte Architektur die Replikation und horizontale Skalierung von Komponenteninstanzen unterstützt, sodass das System bei wachsendem Anrufaufkommen durch Hinzufügen weiterer Container-Replicas skalierbar bleibt.

\textit{Kontext:} Containerorchestrierungsplattformen ermöglichen eine automatisierte Bereitstellung, Scaling und Loadbalancing. Nach \gls{gls-ETSI}-\gls{gls-NFV}-Standard können Netzwerkfunktionen vertikal und horizontal skaliert werden, um dynamisch auf sich ändernde Lasten zu reagieren. Dies wird durch eine dezentrale Komponentenarchitektur und zustandslose Services realisiert \parencite{ETSI-NFV,ETSI-REL,Antonova2025}.

\paragraph{H5 (STT-Robustheit):} Es wird angenommen, dass die eingesetzte Spracherkennungskomponente deutschsprachige Telefonie-Prompts und typische Anrufervorgaben zuverlässig transkribiert, ohne dabei systematische oder regelmä\ss ig wiederholte Erkennungsfehler zu erzeugen.

\textit{Kontext:} Moderne Spracherkennungsmodelle wie OpenAI Whisper sind auf multilingualen Trainingsdaten trainiert und zeigen besonders für Deutsch eine hohe Erkennungsgenauigkeit. Dies stützt die Annahme, dass für telefonische Anwendungen mit kontrollierten Prompt-Szenarien eine zuverlässige Performance zu erwarten ist, auch wenn Test-Bedingungen und reale Produktionsumgebungen unterschiedliche Fehlerraten aufweisen können \parencite{whisper_quant_2025}.
